% CLAUDE.md working-context file, converted to LaTeX for the AI logbook appendix.
% This is a FRAGMENT: it has no preamble and is meant to be \input inside the
% appendices, as a \subsection sibling of "Use of Generative AI".
%
% Relies only on packages already loaded by thesis.tex:
%   array  (>{\raggedright\arraybackslash}p{} columns), enumitem (list labels),
%   longtable (page-breaking tables).  No new packages required.
%
% Heading levels here: \subsection (the file) -> \subsubsection (its ## headings)
% -> \paragraph (its ### headings). Note: with tocdepth=4 in thesis.tex, the
% \paragraph headings will show up in the Table of Contents. If you want the
% appendix kept out of that fine detail, star them: \subsubsection* / \paragraph*.

%----------------------------------------------------------------------
\subsubsection{About}

Matteo Migliaccio (1862669), LLM Law and Technology, Utrecht University.

\textbf{Topic:} GDPR cross-border enforcement and the Procedural Regulation
(Reg. 2025/2518), assessed against the general principle of good administration
under EU law.

\textbf{Central argument:} The GDPR's cross-border enforcement suffers from
structural flaws --- chiefly the LSA's unilateral control over investigation
scope, leaving SACs and the EDPB with no evidentiary record to act on.
Reg. 2025/2518 introduces early consensus finding and procedural deadlines. The
thesis assesses whether these mechanisms adequately curb LSA discretionary power
against the duty of care and duty to give reasons under the objective GPL of good
administration.

%----------------------------------------------------------------------
\subsubsection{Research Questions}

\textbf{Main RQ:} To what extent does Regulation 2025/2518 curb LSAs'
discretionary powers in light of the duty of care and duty to give reasons under
the general principle of good administration in the context of composite
cross-border enforcement of the GDPR?

\begin{itemize}
  \item \textbf{SQ1:} What do the duty of care and duty to give reasons require
        of DPAs in cross-border GDPR enforcement?
  \item \textbf{SQ2:} What changes does the Procedural Regulation bring to the
        GDPR's cooperation and consistency mechanism?
  \item \textbf{SQ3:} Are earlier consensus finding and new procedural deadlines
        under Reg. 2025/2518 adequate to fulfil the duty of care of DPAs?
  \item \textbf{SQ4:} Will earlier consensus finding and new procedural deadlines
        contribute to respecting the duty to give reasons in composite
        administrative proceedings involving the EDPB?
  \item \textbf{SQ5:} What consequences stem from a failure to observe the
        principle of good administration for the protection of fundamental rights
        in the context of cross-border enforcement of the GDPR?
\end{itemize}

%----------------------------------------------------------------------
\subsubsection{AI Use Rules (Utrecht University --- Strict)}

\paragraph{Claude MAY}
\begin{itemize}
  \item Diagnose structural gaps, missing steps, weak justifications, missing sources
  \item Suggest structure, outline alternatives, map argument flow
  \item Ask questions that help Matteo identify what to write
  \item Help locate academic sources (Zotero semantic search, boolean strings)
  \item Brainstorm angles and approaches
\end{itemize}

\paragraph{Claude MUST NOT}
\begin{itemize}
  \item Write thesis text, arguments, or literature/case law summaries
  \item Rewrite or paraphrase Matteo's own text
  \item Generate footnotes or source summaries for direct incorporation
  \item Grammar/style checking until after the final draft is submitted and discussed
\end{itemize}

\textbf{Matteo writes everything himself. Claude provides pointers, not prose.}

\paragraph{AI log}
Remind Matteo to update his AI interaction log at the end of every substantive
session. Log: \texttt{./AI/}

%----------------------------------------------------------------------
\subsubsection{Deadlines}

\begin{longtable}{@{}>{\raggedright\arraybackslash}p{0.68\textwidth}
                    >{\raggedright\arraybackslash}p{0.22\textwidth}@{}}
  \hline
  \textbf{Milestone} & \textbf{Date} \\
  \hline
  \endhead
  Full thesis draft (complete with conclusions, footnotes, bibliography, ToC)
    & \textbf{1 June 2026} \\
  Final thesis submission (PDF via Brightspace)
    & \textbf{26 June 2026} \\
  Reflection report
    & \textbf{1 July 2026} \\
  \hline
\end{longtable}

%----------------------------------------------------------------------
\subsubsection{Requirements}

\begin{itemize}
  \item \textbf{Word count:} 12,500--17,500 (footnotes included;
        bibliography / ToC / title page / abbreviations / annexes excluded).
        Target $\sim$15,000.
  \item \textbf{Citation style:} OSCOLA throughout
  \item \textbf{Font:} Calibri or TNR, 12pt, justified, 1.3--1.5 line spacing,
        space between paragraphs
  \item \textbf{Cover page:} programme, title, name + student number, word count,
        supervisor + second reader, submission date
  \item \textbf{Submission:} PDF via Brightspace
  \item \textbf{Versioning:} Keep multiple versions of text (required under
        Utrecht AI policy)
\end{itemize}

%----------------------------------------------------------------------
\subsubsection{File Locations}

\begin{longtable}{@{}>{\raggedright\arraybackslash}p{0.30\textwidth}
                    >{\raggedright\arraybackslash}p{0.62\textwidth}@{}}
  \hline
  \textbf{Thing} & \textbf{Path} \\
  \hline
  \endhead
  \textbf{This folder} (thesis root, LaTeX project, git-tracked)
    & \texttt{\textasciitilde/Documents/thesis/} locally;
      \texttt{/home/user/thesis/} in Claude Code on the web. Relative paths below
      are unaffected. \\
  GitHub remote
    & configured on \texttt{origin/main} \\
  Main LaTeX entrypoint
    & \texttt{./thesis.tex} \\
  Title page
    & \texttt{./titlepage.tex} \\
  Chapter files
    & \texttt{./Chapters/chapter\_1.tex} \ldots{} \texttt{chapter\_7.tex} \\
  Acronyms / glossary entries
    & \texttt{./abbreviations.tex} --- loaded by \texttt{glossaries} via
      \texttt{\textbackslash{}loadglsentries\{abbreviations\}} in
      \texttt{thesis.tex} \\
  Bibliography
    & \texttt{./bibliography/zotero.bib} (auto-generated by Zotero) +
      \texttt{./bibliography/law.bib} (manual entries) --- OSCOLA via biblatex \\
  latexmk config
    & \texttt{./.latexmkrc} --- pins xelatex, sends artifacts to
      \texttt{Output/}, regenerates glossary/index aux files \\
  Compiled output (PDF, aux files)
    & \texttt{./Output/} (latexmk build dir, gitignored) \\
  Figures
    & \texttt{./Figures/} \\
  Proposal folder (gitignored, reference only)
    & \texttt{./Proposal/} \\
  Submitted proposal
    & \texttt{./Proposal/Thesis Proposal -- Migliaccio.docx} \\
  Proposal draft + comments
    & \texttt{./Proposal/Draft/} \\
  Supervisor feedback (gitignored)
    & \texttt{./Proposal/supervisor-feedback-2026-03-13.md} \\
  AI session logs (gitignored)
    & \texttt{./AI/} \\
  \textbf{Obsidian thesis notes}
    & \texttt{\textasciitilde/Library/Mobile Documents/iCloud\textasciitilde{}md\textasciitilde{}obsidian/Documents/Matteo's Notebook/00 - Thesis Trajectory/} \\
  Introduction template
    & \texttt{.../Proposal/Introduction template.md} (Obsidian) \\
  Literature reading notes
    & \texttt{.../Literature/Notes/} (Obsidian) \\
  Admin docs
    & \texttt{.../Admin/} (Obsidian) \\
  \hline
\end{longtable}

The gitignored locals (\texttt{Proposal/}, \texttt{AI/}, \texttt{flowchart/},
\texttt{Output/}, \texttt{.claude/}) and the Obsidian paths only exist on
Matteo's machine --- they are absent in Claude Code on the web sessions, so don't
assume their contents are readable from the cloud.

%----------------------------------------------------------------------
\subsubsection{Document Structure}

Each chapter: \textbf{intro preview $\rightarrow$ doctrine $\rightarrow$ analysis
$\rightarrow$ conclusion feeding into the next chapter.}

\textbf{Introduction (chapter 1) --- current structure:}

\begin{enumerate}
  \item \textit{Background, Problem, and Analytical Framework} --- integrates lit
        review by substantive point. Subsubsections: Enforcement Challenges of the
        GDPR; The Procedural Regulation; Peer Pressuring the LSA into Action?; The
        Principle of Good Administration as a Benchmark; The Duty of Care; The Duty
        to Give Reasons
  \item \textit{Research Gap}
  \item \textit{Research Questions} (main RQ + 5 sub-RQs)
  \item \textit{Methodology} --- Methods; Limitations; Use of Generative AI
  \item \textit{Outline}
\end{enumerate}

\textbf{General chapters:} each linked to a sub-question. Opens with what the
chapter will do, closes with what it showed and how it feeds into the RQ.

\textbf{Conclusion:} restate problem + RQ $\rightarrow$ summarize chapter findings
$\rightarrow$ state answer $\rightarrow$ limitations + future research.

%----------------------------------------------------------------------
\subsubsection{Key Scholars \& Concepts}

\textbf{Scholars:} Lynskey \& Gentile (GDPR enforcement flaws), Paul Craig (three
functions of reason-giving, 2012), Evans (right to good administration), Mihaescu,
Mustert (supervisor)

\textbf{Institutions:} LSA, SAC, EDPB, DPAs, CJEU, EDPS

\textbf{Concepts:} composite administrative proceedings, OSS mechanism,
consistency mechanism, duty of care, duty to give reasons, GPL of good
administration, Regulation 2025/2518, principle of proximity, early consensus
finding

\textbf{Zotero search tips:} Use semantic search first. Good queries: ``lead
supervisory authority discretion'', ``EDPB composite procedure'', ``duty to give
reasons DPA'', ``good administration GDPR enforcement'', ``principle of proximity
data protection'', ``Paul Craig reason giving''

%----------------------------------------------------------------------
\subsubsection{LaTeX \& Git Conventions}

\begin{itemize}
  \item Thesis is written in LaTeX, edited in VS Code. Draft lives in
        \texttt{Chapters/chapter\_*.tex}, compiled via \texttt{thesis.tex}.
  \item Before suggesting LaTeX commands, check \texttt{thesis.tex} preamble for
        loaded packages --- don't assume what's available.
  \item When introducing an acronym, define it once in \texttt{abbreviations.tex}
        and reference it in prose with \texttt{\textbackslash{}gls\{key\}}.
        Existing keys: \texttt{gdpr}, \texttt{eu}, \texttt{ms}, \texttt{dpa},
        \texttt{oss}, \texttt{lsa}, \texttt{sa}, \texttt{sac}, \texttt{edpb},
        \texttt{edps}, \texttt{cfr}, \texttt{tfeu}, \texttt{cjeu}, \texttt{mcp}.
  \item \texttt{.gitignore} excludes \texttt{.DS\_Store}, \texttt{Proposal/},
        \texttt{AI/}, \texttt{flowchart}, \texttt{Output/}, \texttt{.claude/}, and
        \texttt{*.code-workspace}.
  \item Do not commit on Matteo's behalf unless asked. When asked, use clear
        present-tense messages (e.g., \texttt{add duty of care scaffold to
        chapter 2}).
  \item Never run \texttt{git push -{}-force}, branch deletions, or destructive
        resets without explicit confirmation.
\end{itemize}

\paragraph{Build \& compilation}
\begin{itemize}
  \item \textbf{Command:} \texttt{latexmk thesis.tex} from the repo root.
        \texttt{.latexmkrc} selects \texttt{xelatex} (required by \texttt{fontspec}
        + the Scala Pro mainfont) and writes artifacts to \texttt{Output/}.
  \item \textbf{Glossary + indexes:} \texttt{thesis.tex} loads \texttt{glossaries}
        (acronyms from \texttt{abbreviations.tex}) and \texttt{imakeidx} with two
        indexes --- \texttt{cases} and \texttt{legislation} --- wired into
        biblatex-oscola via \texttt{indexing=cite}. \texttt{.latexmkrc} adds custom
        dependencies so latexmk regenerates \texttt{.gls}/\texttt{.acr} and rebuilds
        the indexes automatically; a single \texttt{latexmk} run is enough.
  \item \textbf{Bibliography:} biblatex with \texttt{style=oscola},
        \texttt{ibid=true}. \texttt{bibliography/zotero.bib}. New manual legal
        entries (cases, treaties, regulations, directives, decisions) go in
        \texttt{bibliography/law.bib}.
  \item \textbf{Cloud caveat:} xelatex and the Scala Pro font may not be installed
        in Claude Code on the web sessions. Don't claim a successful build without
        actually running it --- flag missing toolchain explicitly.
\end{itemize}

\paragraph{LaTeX skill level: beginner}
Matteo is new to LaTeX. When touching \texttt{.tex} files or suggesting LaTeX work:
\begin{itemize}
  \item Briefly explain what commands/environments/packages do --- don't assume
        familiarity
  \item Point out which file change affects what in the compiled PDF
  \item Prefer simple, readable constructs over clever macros
  \item Flag when a package needs to be added to \texttt{thesis.tex} preamble
  \item If a build error appears, explain the error message in plain terms before
        proposing a fix
  \item Remember the AI Use Rules still apply: Claude explains LaTeX mechanics,
        Matteo writes the thesis prose
\end{itemize}

%----------------------------------------------------------------------
\subsubsection{Session Startup}

Do this automatically --- never ask Matteo to re-explain the thesis.

\begin{enumerate}
  \item Read the current draft directly from \texttt{Chapters/chapter\_*.tex} ---
        \textbf{always re-read}, never rely on cached versions. Skim
        \texttt{thesis.tex} too: it's the source of truth for which packages and
        macros are available, so don't suggest a new package without checking the
        preamble first.
  \item If doing source work: check Zotero semantic search database status, update
        if needed
  \item Check Things 3 for thesis deadlines if relevant
  \item Check Obsidian notes (\texttt{00 - Thesis Trajectory/}) if relevant
\end{enumerate}

\paragraph{For structural review}
\begin{enumerate}
  \item Read current \texttt{Chapters/chapter\_*.tex} files
  \item Read the introduction template from Obsidian
  \item Compare against supervisor feedback points above
  \item Identify gaps --- do NOT write the text
\end{enumerate}

\paragraph{For source finding}
\begin{enumerate}
  \item \texttt{zotero\_semantic\_search} first (topic discovery)
  \item \texttt{zotero\_search\_items} for known authors (keep queries short:
        ``Lynskey 2024'', not long phrases)
  \item \texttt{zotero\_get\_annotations} to check what Matteo already highlighted
  \item \texttt{zotero\_get\_pdf\_outline} before \texttt{zotero\_get\_item\_fulltext}
        to scope what to read
  \item \textbf{Context warning:} \texttt{zotero\_get\_item\_fulltext} returns 10K+
        tokens. Avoid on multiple papers per session.
\end{enumerate}

%----------------------------------------------------------------------
\subsubsection{End of Session}

\begin{itemize}[label={\framebox[1.4ex]{\rule{0pt}{1.1ex}}}]
  \item Remind Matteo to log this AI interaction (date, tool, prompt, purpose,
        reflection)
  \item Update Things 3 if new deadlines or tasks emerged
  \item Save structural plans to Obsidian if useful for next session
\end{itemize}
